%% ============================================================================
%% Section III: Corpus and Data Preprocessing
%% ============================================================================

\section{Corpus and Data Preprocessing}
\label{sec:preprocessing}

The offline pipeline (Fig.~\ref{fig:offline}) turns 66 books of the CUV
from PDF into (i) a hierarchical structured corpus, (ii) a multi-granular
embedding queue, and (iii) the inputs of knowledge-graph construction. All
build steps run on the host from a dedicated \code{scripts/} environment;
every intermediate artifact is materialized as JSONL before database
import, so each stage is independently inspectable and re-runnable.

\begin{figure*}[!t]
\centering
\resizebox{0.98\textwidth}{!}{%
\begin{tikzpicture}[
  font=\sffamily\footnotesize,
  >={Stealth[length=2mm,width=1.6mm]},
  box/.style={rectangle, rounded corners=2pt, draw=black!60, fill=black!6,
              align=center, inner sep=3pt, minimum height=0.85cm},
  data/.style={rectangle, rounded corners=2pt, draw=blue!55, fill=blue!8,
               align=center, inner sep=3pt, minimum height=0.85cm},
  db/.style={cylinder, shape border rotate=90, aspect=0.18, draw=teal!65,
             fill=teal!12, align=center, inner sep=2.5pt},
  flow/.style={->, draw=black!62, line width=0.7pt},
]
% Row 1: source to chunking
\node[data] (pdf) {66-book PDF\\(CUV)};
\node[box, right=0.55cm of pdf] (md) {Layout-geometry\\PDF$\to$Markdown};
\node[box, right=0.55cm of md] (parse) {Hierarchical parsing\\+ chunking};
\node[data, right=0.55cm of parse, text width=2.9cm] (queue)
  {Embedding queue\\34{,}072 units\\(pericope/chunk/verse)};
\node[data, right=0.55cm of queue, text width=2.6cm] (struct)
  {Structure JSONL\\books/chapters/\\pericopes/chunks\\+ cross-refs};

% Row 2: extraction / embedding
\node[box, below=0.8cm of queue.south west, anchor=north west, text width=2.9cm] (ee)
  {Entity extraction\\(dual-track, grounded)};
\node[box, right=0.5cm of ee, text width=2.9cm] (re)
  {Relation extraction\\(37-type closed schema)};
\node[box, left=0.5cm of ee, text width=2.6cm] (emb)
  {BGE-M3 dense\\+ BM25/CKIP sparse\\+ entity vectors};
\node[data, right=0.5cm of re, text width=2.2cm] (tsk)
  {TSK cross-refs\\(public domain)};

% Row 3: databases
\node[db, below=0.9cm of emb, minimum width=2.3cm, minimum height=1.05cm] (qd)
  {\textbf{Qdrant}\\3 collections};
\node[db, below=0.9cm of ee, minimum width=2.3cm, minimum height=1.05cm] (pg)
  {\textbf{PostgreSQL}\\6 tables};
\node[db, below=0.9cm of re, minimum width=2.5cm, minimum height=1.05cm] (neo)
  {\textbf{Neo4j}\\13{,}589 nodes\\319{,}988 edges};

\draw[flow] (pdf) -- (md);
\draw[flow] (md) -- (parse);
\draw[flow] (parse) -- (queue);
\draw[flow] (parse.north) -- ++(0,0.28) -| (struct.north);
\draw[flow] (queue.south) |- ($(ee.north)+(0,0.28)$) -- (ee.north);
\draw[flow] (queue.south) |- ($(emb.north)+(0,0.28)$) -- (emb.north);
\draw[flow] (ee) -- (re);
\draw[flow] (emb) -- (qd);
\draw[flow] (ee) -- (pg);
\draw[flow] (re) -- (neo);
\draw[flow] (tsk.south) |- ($(neo.east)+(0.25,0.15)$) -- ($(neo.east)+(0,0.15)$);
\draw[flow] (struct.east) -| ++(0.35,-3.1) |- ($(neo.east)+(0,-0.15)$);
\draw[flow] (ee.south) -- ++(0,-0.33) -| ($(neo.north)+(-0.35,0)$);
\draw[flow] (struct.south) -- ++(0,-0.4) -| ($(pg.north)+(0.4,0)$);
\end{tikzpicture}%
}
\caption{Offline build pipeline. Every arrow is a materialized JSONL
hand-off; the three databases play complementary roles
(Section~\ref{sec:tripledb}); database counts show the live
post-remediation state (July 6, 2026). The embedding queue doubles as
the extraction feed --- the structural root cause analyzed in
Section~\ref{sec:audit} (mechanism M1).}
\label{fig:offline}
\end{figure*}

\subsection{Source Corpus}

The source is the Chinese Union Version with new punctuation, obtained as
66 per-book PDFs (one per biblical book). The CUV is the de facto standard
Protestant translation in the Traditional Chinese world; its vocabulary is
early-20th-century literary Chinese, which creates a systematic lexical
gap between modern user phrasing and corpus surface forms --- a gap that
recurs throughout this paper (\eg the modern phrase ``the Last Supper''
never occurs in the CUV text, which narrates a ``Passover meal'').

\subsection{PDF-to-Markdown by Layout Geometry}

Conversion (\code{convert\_bible\_pdf.py}, PyMuPDF) uses \emph{pure layout
geometry} over character-level spans --- no OCR and no LLM --- so the step
is deterministic, free, and auditable:

\begin{itemize}
\item font size $\le 8.5$\,pt: running headers/footers, dropped;
\item font size $\ge 16$\,pt: book title or chapter number;
\item font size $=12$\,pt with $x \ge 85$: pericope (section) heading;
\item font size $\le 9.5$\,pt and numeric: verse number;
\item verse number at $x < 75$ with line right-edge $< 300$: poetry
      (line breaks preserved); otherwise prose (lines joined);
\item bold \code{N:N} tokens: footnotes, aggregated per chapter.
\end{itemize}

The resulting Markdown uses \code{\#} for book, \code{\#\#} for chapter,
\code{\#\#\#} for pericope headings, and bold verse numbers. Under each
pericope heading, the CUV's printed parallel-passage annotations (\eg
``(Mark 1:9--11; Luke 3:21--22)'' under the baptism pericope) are retained
--- they become the first source of the cross-reference network
(Section~\ref{sec:crossrefs}).

\subsection{Pericope-Centric Hierarchical Structure}
\label{sec:pericope}

A \emph{pericope} is the established interpretive unit of biblical
scholarship: a self-contained narrative episode, discourse, or psalm. The
CUV's editorial pericope headings give us this segmentation for free, and
the system adopts the pericope as the core unit for embedding, entity
anchoring, and cross-referencing, for three reasons:

\begin{enumerate}
\item \textbf{Semantic completeness.} A story is never split down the
middle; retrieval returns interpretable wholes.
\item \textbf{Embedding fit.} The median pericope is 388 characters
(P90 $=931$, max $=3{,}544$), well inside the effective input range of the
BGE-M3 embedder~\cite{chen2024bgem3}.
\item \textbf{Explainability.} Returning ``Isaac takes a wife
(Gen.~24)'' is more interpretable to end users than a verse-range span.
\end{enumerate}

Parsing (\code{process\_bible.py} with \code{markdown\_parser.py}) yields
the hierarchy Book\,$\rightarrow$\,Chapter\,$\rightarrow$\,Pericope with
counts 66\,/\,1{,}189\,/\,2{,}779, plus per-pericope verse spans. A
regex-based reference parser with a Chinese book-abbreviation table
resolves the parallel-passage annotations; every reference is
\emph{upgraded to pericope level} (never verse level), keeping the
cross-reference network aligned with the retrieval unit. Hierarchical IDs
are colon-delimited (\code{gen} $\rightarrow$ \code{gen:1} $\rightarrow$
\code{gen:1:0} for the first pericope of Genesis~1), a scheme reused
verbatim across all three databases.

\subsection{Hierarchical Chunking}
\label{sec:chunking}

Long pericopes are split for embedding by a token-aware hierarchical
chunker whose tokenizer is \emph{BGE-M3's own}, i.e., the chunking layer
is explicitly designed for the embedding model~\cite{chen2024bgem3}:

\begin{itemize}
\item pericopes with $\le 768$ tokens (93.9\%) are kept whole;
\item longer pericopes are split on \emph{verse boundaries} targeting 512
tokens per chunk, with a 1-verse overlap between adjacent chunks;
\item a trailing chunk under 128 tokens and $\le 3$ verses is merged back
into its predecessor;
\item chunk IDs extend the parent ID (\code{gen:1:0:0}), preserving the
hierarchy.
\end{itemize}

Only 169 of 2{,}779 pericopes (6.1\%) are split, producing 431 chunks; the
split pericopes concentrate in Old Testament narrative books (Leviticus
13, Judges 12, 1~Samuel 12, Genesis 11, Numbers 11, Joshua 10, Deuteronomy
9, 2~Samuel 9) and account for 18.3\% of the corpus text (243{,}773 of
1{,}329{,}046 characters); their median length is 1{,}353 characters
(max 3{,}544), against a chunk median of 644 (max 796). Each retrieval
unit's embedded text carries a locating prefix ---
``\emph{Book, Chapter $N$, pericope title (verses $X$--$Y$):}'' --- so
vectors encode provenance context. Table~\ref{tab:corpus} summarizes the
corpus.

\begin{table}[!t]
\caption{Corpus and Chunking Statistics (CUV, New Punctuation)}
\label{tab:corpus}
\centering
\begin{tabular}{lr}
\toprule
Item & Value \\
\midrule
Books & 66 \\
Chapters & 1{,}189 \\
Pericopes & 2{,}779 \\
\quad -- split into chunks & 169 (6.1\%) \\
Chunks (from split pericopes) & 431 \\
Corpus size (characters) & 1{,}329{,}046 \\
\quad -- inside split pericopes & 243{,}773 (18.3\%) \\
Pericope length, median / P90 / max (chars) & 388 / 931 / 3{,}544 \\
Split-pericope length, median / max (chars) & 1{,}353 / 3{,}544 \\
Chunk length, median / max (chars) & 644 / 796 \\
\midrule
Chunker: split threshold / target / min (tokens) & 768 / 512 / 128 \\
Chunker: overlap & 1 verse \\
Tokenizer & BGE-M3 \\
\bottomrule
\end{tabular}
\end{table}

\subsection{Multi-Granularity Embedding Queue}
\label{sec:queue}

Retrieval needs different granularities for different question shapes, so
the pipeline emits an embedding queue that is the union of three levels,
34{,}072 units in total:

\begin{itemize}
\item \textbf{pericope level} (2{,}610): every pericope that was
\emph{not} split;
\item \textbf{chunk level} (431): split pericopes enter only as chunks;
\item \textbf{verse level} (31{,}031): every verse individually, for
fine-grained exact retrieval (printed verse ranges such as ``1--2'' are
kept as single units).
\end{itemize}

This queue is an artifact designed \emph{for retrieval}. Critically for
what follows, the early entity-extraction pipeline consumed this same
queue as its input feed. That single wiring decision --- extraction feed
$=$ embedding queue --- silently propagated an embedding-oriented
granularity into the knowledge graph and became the structural root cause
(M1) of the defect chain quantified in Section~\ref{sec:audit}.
